% =========================================================
% config/listings.tex
% Konfigurasi penampilan source code seluruh buku.
% Bahasa utama: assembly 8086 dengan sintaks MASM/TASM.
% =========================================================

% ---------------------------------------------------------
% 1. Definisi bahasa: MASM/TASM 8086 (case-insensitive)
% ---------------------------------------------------------
\lstdefinelanguage{MASM}{
    morekeywords={
        % --- transfer data ---
        MOV, XCHG, LEA, LDS, LES, PUSH, POP, PUSHA, POPA, PUSHF, POPF,
        LAHF, SAHF, XLAT, IN, OUT, CBW, CWD,
        % --- aritmatika ---
        ADD, ADC, SUB, SBB, INC, DEC, NEG, CMP, MUL, IMUL, DIV, IDIV,
        AAA, AAS, AAM, AAD, DAA, DAS,
        % --- logika dan geser ---
        AND, OR, XOR, NOT, TEST, SHL, SHR, SAL, SAR, ROL, ROR, RCL, RCR,
        % --- kontrol alur ---
        JMP, JE, JZ, JNE, JNZ, JG, JGE, JL, JLE, JA, JAE, JB, JBE,
        JC, JNC, JS, JNS, JO, JNO, JP, JPE, JNP, JPO, JCXZ,
        LOOP, LOOPE, LOOPZ, LOOPNE, LOOPNZ, CALL, RET, INT, IRET, INTO,
        % --- prosesor ---
        HLT, NOP, CLI, STI, CLD, STD, CLC, STC, CMC, WAIT, ESC, LOCK,
        % --- direktif data ---
        DB, DW, DD, DQ, DT, EQU, ORG, SEGMENT, ENDS, ASSUME, END,
        PROC, ENDP, MACRO, ENDM, LOCAL, PUBLIC, EXTRN, EXTERN,
        INCLUDE, INCLUDELIB, MODEL, STACK, CODE, DATA, CONST, ALIGN, EVEN,
        % --- operator ---
        OFFSET, PTR, BYTE, WORD, DWORD, QWORD, TBYTE,
        NEAR, FAR, SHORT, LABEL, TYPE, SIZE, LENGTH, DUP, THIS,
        % --- register ---
        AX, BX, CX, DX, AL, AH, BL, BH, CL, CH, DL, DH,
        SP, BP, SI, DI, IP, CS, DS, SS, ES, FLAGS,
        % --- flag ---
        CF, PF, AF, ZF, SF, TF, IF, DF, OF,
        % --- pseudo-op umum ---
        BYTE, NOTHING, TRUE, FALSE
    },
    sensitive=false,
    morecomment=[l]{;},
    morestring=[b]',
    morestring=[b]"
}

% ---------------------------------------------------------
% 2. Gaya dasar
% ---------------------------------------------------------
\lstdefinestyle{asmgaya}{
    language=MASM,
    basicstyle=\ttfamily\footnotesize,
    keywordstyle=\color{obeBlue}\bfseries,
    commentstyle=\color{obeGreen}\itshape,
    stringstyle=\color{obeMaroon},
    numberstyle=\tiny\color{obeGray},
    numbers=left,
    numbersep=8pt,
    stepnumber=1,
    backgroundcolor=\color{codeBack},
    frame=single,
    rulecolor=\color{codeFrame},
    framesep=6pt,
    breaklines=true,
    breakatwhitespace=false,
    postbreak=\mbox{\textcolor{obeGray}{$\hookrightarrow$}\space},
    showstringspaces=false,
    tabsize=4,
    columns=fullflexible,
    keepspaces=true,
    upquote=true,
    xleftmargin=14pt,
    aboveskip=0.8em,
    belowskip=0.6em,
    captionpos=b
}

\lstset{style=asmgaya}

% ---------------------------------------------------------
% 3. Environment listing bernomor untuk program utuh
% ---------------------------------------------------------
\lstnewenvironment{program}[1][]{%
    \lstset{style=asmgaya, #1}%
}{}

% ---------------------------------------------------------
% 4. Gaya untuk keluaran program / sesi debugger (tanpa sorot)
% ---------------------------------------------------------
\lstdefinestyle{sessiongaya}{
    basicstyle=\ttfamily\footnotesize,
    keywordstyle=\color{black},
    commentstyle=\color{obeGray},
    stringstyle=\color{black},
    numbers=none,
    backgroundcolor=\color{obeGrayLight},
    frame=single,
    rulecolor=\color{codeFrame},
    framesep=6pt,
    breaklines=true,
    showstringspaces=false,
    columns=fullflexible,
    keepspaces=true,
    upquote=true,
    xleftmargin=6pt,
    aboveskip=0.8em,
    belowskip=0.6em
}

\lstnewenvironment{session}[1][]{%
    \lstset{style=sessiongaya, #1}%
}{}

% ---------------------------------------------------------
% 5. Gaya untuk pseudocode
% ---------------------------------------------------------
\lstdefinestyle{pseudogaya}{
    basicstyle=\ttfamily\footnotesize,
    keywordstyle=\color{obeTeal}\bfseries,
    numbers=left,
    numberstyle=\tiny\color{obeGray},
    backgroundcolor=\color{obeGrayLight},
    frame=single,
    rulecolor=\color{codeFrame},
    framesep=6pt,
    breaklines=true,
    showstringspaces=false,
    columns=fullflexible,
    keepspaces=true,
    upquote=true,
    xleftmargin=14pt
}

\lstnewenvironment{pseudocode}[1][]{%
    \lstset{style=pseudogaya, #1}%
}{}

% ---------------------------------------------------------
% 6. \lstinline dengan gaya seragam
% ---------------------------------------------------------
\newcommand{\lstasm}{%
    \lstinline[style=asmgaya,numbers=none,basicstyle=\ttfamily\small,frame=none,backgroundcolor=\color{white}]%
}
